\subsection{C4：规则化清洗的典型代表}

C4 是 T5 论文里最重要的数据贡献之一。它使用大量手工规则清洗 Common Crawl，比如只保留以标点结尾、且足够长的行，过滤少句页面、脏词、占位符和 boilerplate，并只保留英语。具体规则包括：

\begin{itemize}
    \item 只保留以标点符号结尾且包含 $\geq 5$ 个单词的行。
    \item 删除少于 3 个句子的页面。
    \item 删除包含任何"脏词"的页面（基于公开的脏词列表）。
    \item 删除包含 \texttt{\{} 的页面（排除代码）、包含 \texttt{lorem ipsum} 的页面、包含 \texttt{terms of use} 的页面等。
    \item 使用 langdetect 过滤，只保留英语概率 $> 0.99$ 的页面。
\end{itemize}

起点是 2019 年 4 月的 Common Crawl 快照（约 1.4T tokens），经过上述规则过滤后得到 806GB 文本（约 156B tokens）。

\begin{warningbox}{规则清洗的代价}
规则简单意味着透明，但也意味着误杀和漏网都很多。它能保留一些不太像 Wikipedia、但很有用的句子，也会放进一些语义上并不理想的内容。例如，"包含 \texttt{\{} 就删除"这条规则会误杀所有包含 JSON 或数学公式的页面，而"脏词列表"则可能把合法的医学或性教育内容一并排除。
\end{warningbox}

一个有趣的附加实验是：T5 团队还做了一个 WebText-like 的子集，只保留出现在 OpenWebText 链接中的页面（即 Reddit karma $\geq 3$ 的外链）。用了 12 个 Common Crawl dump 才凑到 17GB——而原版 WebText 有 40GB，这说明 Common Crawl 的覆盖面虽广，但远非完整。

\subsubsection{Worked Example：去重方法对比}

去重是数据清洗中最关键的步骤之一。不同去重方法在精度、召回率和计算成本上差异显著：

\begin{table}[H]
\centering
\begin{tabular}{p{0.15\textwidth}p{0.22\textwidth}p{0.2\textwidth}p{0.18\textwidth}p{0.15\textwidth}}
\toprule
\textbf{方法} & \textbf{原理} & \textbf{优点} & \textbf{缺点} & \textbf{典型用户} \\
\midrule
精确去重 & 计算文档或段落的 SHA-256 哈希，完全匹配则删除 & 零误杀，实现简单 & 对细微修改无效（改一个字就不匹配） & CCNet \\
MinHash + LSH & 用 n-gram 的最小哈希估计 Jaccard 相似度，将相似文档聚在同一桶内 & 能捕捉近重复（如换了标题或广告的相同文章） & 有一定误杀率；超参选择影响大 & RefinedWeb, FineWeb, SlimPajama \\
Bloom filter & 用概率数据结构记录已见过的 n-gram 或 URL & 内存高效 & 有假阳性 & Dolma \\
Suffix array & 找到精确重复的长子串 & 可以做子文档级去重 & 计算成本高 & The Pile \\
\bottomrule
\end{tabular}
\caption{去重方法对比：没有银弹，需要根据数据规模和质量目标选择}
\end{table}

\begin{warningbox}{数据去重不充分会导致模型记忆}
研究已经证明，训练数据中的重复内容会导致模型\textbf{逐字记忆}（verbatim memorization）。具体表现为：给模型一段重复出现过的文本的前缀，它能几乎完美地续写出后续内容。这不仅是隐私风险（如果训练数据包含个人信息），也是版权风险（如果模型能复述受版权保护的作品）。Carlini et al.~(2023) 的研究表明，一段文本在训练集中出现的次数越多，模型能完整复述它的概率就越高。去重不是"锦上添花"，而是防止模型变成搜索引擎的基本功。
\end{warningbox}

